\subsection{完全切割序列的一个判据}

设 A 为一个环，M 为一个 A-模，x 为 A 的一个理想。M 上的 x-进拓扑是与 M 的群结构相容的拓扑，且在该拓扑下，子模 \(x^{r}\) M ( \(r \geqslant 0\) ) 的集合构成零的一个基本邻域系（TG，III，第 5 页，例）。该拓扑是分离的当且仅当

\[
\bigcap_ {r \geqslant 0} x ^ {r} = 0.
\]

现在假设理想 x 由 A 中元素的序列 \(\boldsymbol{x}=(x_{1},\ldots,x_{n})\)

生成。考虑分次 A-模 \(\bigoplus_{r\geqslant 0}x^{r}\) M 以及 0 次的分次 A-同态

\[
a _ {\mathbf {M}} ^ {\mathbf {x}}: \mathrm{A} [ \mathrm{X} _ {1}, \dots , \mathrm{X} _ {n} ] \otimes_ {\mathbf {A}} \mathrm{M} \rightarrow \bigoplus_ {r \geqslant 0} \mathfrak {x} ^ {r} \mathrm{M}
\]

使得 \(a_{\mathbf{M}}^{\star}(\mathbf{P} \otimes m) = \mathbf{P}(x_1, \ldots, x_n) \, m\) ，如果 \(\mathbf{P}\) 是 \(\mathbf{X}_1, \ldots, \mathbf{X}_n\) 中的齐次多项式且 \(m \in \mathbf{M}\) 。设 \(\mathfrak{d}\) 为 \(\mathbf{A}[\mathbf{X}_1, \ldots, \mathbf{X}_n]\) 中由元素 \((x_i \, \mathbf{X}_j - x_j \, \mathbf{X}_i)\) （ \(1 \leqslant i < j \leqslant n\) ）生成的理想。我们有 \(\mathbf{P}(x_1, \ldots, x_n) = 0\) ，如果 \(\mathbf{P} \in \mathfrak{d}\) ，从而 \(a_{\mathbf{M}}^{\star}\) 通过过渡到商给出一个 0 次的分次 A-同态

\[
\alpha_ {\mathrm{M}} ^ {x}: \left(\mathrm{A} \left[ \mathrm{X} _ {1}, \dots , \mathrm{X} _ {n} \right] / \mathfrak {d}\right) \otimes_ {\mathrm{A}} \mathrm{M} \rightarrow \bigoplus_ {r \geqslant 0} x ^ {r} \mathrm{M}.
\]

通过与 A/x 作张量积，我们从 \(\alpha_{\mathrm{M}}^{\mathbf{x}}\) 导出一个 0 次的分次 A-同态

\[
\beta_ {\mathrm{M}} ^ {x}: (\mathrm{A} / \mathfrak {x}) [ \mathrm{X} _ {1}, \dots , \mathrm{X} _ {n} ] \otimes_ {\mathrm{A}} \mathrm{M} \rightarrow \bigoplus_ {r \geqslant 0} (\mathfrak {x} ^ {r} \mathrm{M} / \mathfrak {x} ^ {r + 1} \mathrm{M}).
\]

同态 \(a_{\mathbf{M}}^{x}\) 、\(\alpha_{\mathbf{M}}^{x}\) 与 \(\beta_{\mathbf{M}}^{x}\) 是满射的。

定理 1. --- 考虑以下条件：

\begin{enumerate}
\def\labelenumi{(\roman{enumi})}
\item
  序列 x 是 M-正则的（X，第 158 页）。
\item
  序列 x 对 M 完全切割（X，第 157 页，定义 2）。
\item
  我们有 \(\mathbf{H}_1(x, \mathbf{M}) = 0\) 。
\item
  \(L^{\prime}\) 同态 \(\alpha_{\mathbf{M}}^{x}:(\mathbf{A}[\mathbf{X}_{1},\dots,\mathbf{X}_{n}]/\mathfrak{d})\otimes_{\mathbf{A}}\mathbf{M}\to\bigoplus x^{r}\mathbf{M}\) 是双射。
\item
  同态 \(\beta_{\mathbf{M}}^{\mathbf{x}}:(\mathbf{A}/\mathfrak{x})[\mathbf{X}_{1},..., \mathbf{X}_{n}]\otimes_{\mathbf{A}}\mathbf{M}\to\bigoplus_{r\geqslant0}(\mathfrak{x}^{r}\mathbf{M}/\mathfrak{x}^{r+1}\mathbf{M})\) 是双射。
\end{enumerate}

那么我们有蕴含关系 (i)⇒(ii)⇒(iii)⇔(iv)⇒(v)。若对 \(1 \leqslant i \leqslant n\) ，A-模 \(\mathrm{M}/(x_{1}\mathrm{M} + \cdots + x_{i-1}\mathrm{M})\) 对 x-进拓扑是分离的，则条件 (i) 至 (v) 等价。

该定理将在第 8 节和第 9 节中证明。

推论 1. --- 若 A 是诺特环，若 A-模 M 是有限生成的，且若 \(x_{i}\) 属于 A 的根，则定理的条件 (i) 至 (v) 等价。事实上，在每个模 \(\mathbf{M}/(x_{1} \mathbf{M} + \cdots + x_{i-1} \mathbf{M})\) 上 x-进拓扑都是分离的（AC III，§ 3，第 3 号，命题 6）。

推论 2. --- 设 A 是一个正分次环，M 是一个下有界的分次 A-模，且 \(x_i\) 为 A 中次数 \textgreater{} 0 的齐次元素。则定理的条件 (i) 至 (v) 等价。

对任何下有界的分次 A-模 N，x-进拓扑确实是分离的，因为若 \(N_{n}=0\) 对 \(n<n_{0}\) ，则我们有 \(x^{a}N\subset\sum_{j\geqslant n_{0}+a}N_{j}\) 对所有 \(a\geqslant0\) 。

推论 3. --- 假设模 \(\mathbf{M}/(x_1\mathbf{M} + \cdots + x_{i-1}\mathbf{M})\) 对 \(\mathbf{x}\) -进拓扑是分离的（ \(1 \leqslant i \leqslant n\) ）；设 \(p\) 为介于 1 与 \(n\) 之间的整数。序列 \(\mathbf{x}\) 对 \(\mathbf{M}\) 完全切割，当且仅当序列 \((x_1, \ldots, x_p)\) 对 \(\mathbf{M}\) 完全切割且序列 \((x_{p+1}, \ldots, x_n)\) 对 \(\mathbf{M}/(x_1\mathbf{M} + \cdots + x_p\mathbf{M})\) 完全切割。

事实上，如果在陈述中将“完全切割序列”替换为“正则序列”，则该推论是显然的；而根据定理，这两个概念在此处是一致的。

注记。--- 设 \(x = (x_1, \ldots, x_n)\) 是 A 中元素的序列，使得 \(H_1(x, A) = 0\) ；则满射同态 \(u: A^n \to x\) （它使得 \(u(\sum a_i e_i) = \sum a_i x_i\) ）的核由元素 \(X_j e_i - X_i e_j\) 生成；因此，A-代数 \(A[X_1, \ldots, X_n] / \mathfrak{d}\) 同构于对称代数 \(S_A(x)\) （III，第 69 页，命题 4）。于是由定理 1 可知，代数同态 \(S_A(x) \to \bigoplus_n x^n\) （它由 \(x\) 到 \(\bigoplus_n x^n\) 中的典范单射导出）是一个同构。对同态 \(S_A(x / x^2) \to \bigoplus_n x^n / x^{n+1}\) 同样成立。

